\documentclass[../../main2.tex]{subfiles}

\begin{document}

\chapter{通道二：经济（观念 + 稀缺约束）}
\label{ch:economy}

\begin{motivation}
	上一章遗留的问题：模因传播不消耗原体、也就不标价。但现实里大量传播必须付费——印书要钱，买量要钱，办报要钱。如果一种传播\textbf{必须}消耗稀缺资源，它还只是模因吗？\\
	\textbf{核心动机}：用「加入稀缺约束」这一步，从观念通道里把经济\textbf{推导}出来——供求不是经济学的起点，是「观念 + 稀缺」两步走的产物。\\
	\textbf{扩增模式}：有压——复制变转移，总量守恒，零和上场。
\end{motivation}

\begin{logicmap}
\begin{tikzpicture}[node distance=0.85cm, every node/.style={draw, rectangle, rounded corners, align=center, font=\small}]
\node (q) {传播要花钱，还是模因吗？};
\node (s) [below=of q] {加稀缺约束：复制变转移};
\node (sd) [below=of s] {供与求：从摊位到全城菜价};
\node (h) [below=of sd] {预期上场：追涨杀跌};
\node (c) [below=of h] {商路：丝绸之路的账};
\node (n) [below=of c] {交换真的等价吗？（抛给第 13 章）};
\draw[-{Stealth}] (q) -- (s);
\draw[-{Stealth}] (s) -- (sd);
\draw[-{Stealth}] (sd) -- (h);
\draw[-{Stealth}] (h) -- (c);
\draw[-{Stealth}] (c) -- (n);
\end{tikzpicture}
\end{logicmap}

\section{加入稀缺约束：经济从观念里长出来}

\begin{readingnote}
	你有没有想过，为什么「分享」这个词在朋友圈免费，在股市要收印花税？\\
	同一个词，两种压强——差在\textbf{传的东西耗不耗}。
\end{readingnote}

标准句式走一遍。如果传播不消耗原体（你把想法告诉我，你还有这个想法），扩张是无限的——这是上一章的世界。现在加入约束：\textbf{传播需要消耗稀缺资源}——纸张、版面、运力、钱。这一步改变了所有性质：复制变成转移，总量开始守恒，你多我就少。这就是经济行为从文化行为中分离出来的时刻。

\begin{tcolorbox}[colback=green!5, title=\textit{生活类比}]
	上一章是传火，这一章是分粥。火越传越旺，粥越分越薄。经济学全部的紧张感，来自「粥是有限的，火是无限的」，而人类偏偏想用粥的规则去管火，又想用火的规则去煮粥。
\end{tcolorbox}

守恒带来三个新脾气，先挂号，后面几节逐个验：

\begin{itemize}
	\item \textbf{价格}：稀缺的东西必须排队，价格是排队的记分牌；
	\item \textbf{零和}：转移型分配里，我的多是你的少——冲突有了数学根子；
	\item \textbf{预期}：连「未来稀不稀缺」都开始被炒——下一节就到。
\end{itemize}

\paragraph*{逻辑衔接}
	稀缺上场，「分粥」开始。那菜价、房价这些宏观数字，怎么从千万个摊位的算计里冒出来？下一节用你每天都在碰的东西——菜价——走一遍微观到宏观。

\section{从摊位到全城：菜价怎么形成}

\begin{readingnote}
	你有没有想过，菜价每天在变，可\textbf{没人}开过「定价大会」？\\
	千万个摊贩各算各的账，价格却全城联动——谁在记分？
\end{readingnote}

把菜价拆成摊贩的动机，三层就到：

\begin{enumerate}
	\item \textbf{摊主想要什么}：多赚、少压货、比隔壁卖得快。动机不决定结果（第 1 章），他的报价还受环境捏着：今早来货多少、天气、隔壁卖几块。
	\item \textbf{市场怎么汇总}：买菜的人用脚投票——贵了去隔壁，便宜了排队。摊主互相看得见，报价被彼此\textbf{校准}。千万次「贵了不买、贱了抢」，就是价格。
	\item \textbf{宏观规律怎么冒出来}：暴雨断供，全城菜价一起跳——没有合谋，是同一场雨改写了每个摊主的约束。所谓「市场规律」，就是这种「各自算账、互相校准」的统计形状。
\end{enumerate}

\begin{tcolorbox}[colback=green!5, title=\textit{生活类比}]
	价格像河面的水位：每滴雨（每个摊贩的算计）都不想「决定水位」，但水位就是雨滴们互相挤压出来的。看水位能行船（看价格能买卖），想改水位要修河堤（想改价格要动约束——供给、准入、预期）。
\end{tcolorbox}

不同人群的「价格从哪来」现场：

\begin{itemize}
	\item \textbf{买菜的你}：香菜忽然贵三倍——查天气：主产区下雹子了。动机没变，约束变了。
	\item \textbf{学生}：考研资料年年涨价——报名人数是「买菜的人」，名师产能是「来货量」。
	\item \textbf{上班族}：你的工资也是价格——是「劳动力摊位」上的报价，受行业供需捏着，不完全由「能力」定价。
	\item \textbf{古代}：长安粮价随漕运波动——运河就是「来货量」的开关，史书里的粮价曲线跟着河道走。
\end{itemize}

\paragraph*{逻辑衔接}
	静态的价格会看了。但人还会炒「未来」——一旦预期入场，价格会甩开当下的供需，下一节。

\section{预期上场：追涨杀跌}

\begin{readingnote}
	你有没有想过，为什么房价越涨越有人买，菜越跌越没人要？\\
	买的不是房子，是\textbf{「明天更贵」这个念头}。
\end{readingnote}

上一节的价格只算今天。加入新变量：\textbf{预期}。当「未来会涨」本身成了购买理由，需求就自我喂养——这叫追涨杀跌。三个熟悉现场：

\begin{itemize}
	\item \textbf{房价}：「再不买更贵」的念头本身制造购买力，推着「更贵」实现——直到某天念头转向。
	\item \textbf{猪周期}：猪肉涨 $\rightarrow$ 养猪户扩栏 $\rightarrow$ 一年后集中出栏 $\rightarrow$ 肉价崩 $\rightarrow$ 杀猪 $\rightarrow$ 再涨。预期的时滞，画出了那个著名的蛛网。
	\item \textbf{比特币}：没有现金流、没有摊位，纯粹的预期载体——价格几乎全由「下一个人出多少」决定。赌场逻辑穿了件科技外衣。
\end{itemize}

\lowfreq{诚实标注：预期模型（如蛛网模型）是理想化，现实里有大玩家操纵、政策突变、信息不对称——噪声很大，但「预期自我喂养」这个机制反复出现。}

\begin{questionbox}
	\textit{现在你可能想反驳：「追涨杀跌不是『不理性』吗？教科书说人会均值回归。」}\\
	\textbf{个体可能是「理性」的。}「怕上不了车」对每个人都是合理自保，合起来才成了疯——又一个微观动机 $\rightarrow$ 宏观形状的例子（第 10 章那套账）。骂个体不理性没用，要看通道怎么放大。
\end{questionbox}

\paragraph*{逻辑衔接}
	价格与预期都会算了。但经济通道不只在国内转——一条商路能把两个文明的账本接起来。下一节看两千年前的一次「全球化」。

\section{丝绸之路：一条商路的复式账}

\begin{readingnote}
	你有没有想过，一条卖丝绸的路，怎么顺手改了两个文明的\textbf{脑子}？\\
	经济通道跑的不只是货，货里夹着模因。
\end{readingnote}

丝绸之路是经济通道的标本级案例，账要分两边记：

\begin{itemize}
	\item \textbf{中国的账}：丝绸、瓷器换回金银、良马、香料。西汉「凿空」西域，动机是军事联盟（张骞找大月氏），结果长出了商路——又是动机不决定结果。
	\item \textbf{罗马的账}：普林尼抱怨黄金东流，丝绸「比黄金还贵」。奢侈品贸易改写了地中海的消费口味。
	\item \textbf{夹带的模因}：佛教沿商路东传，造纸术西走，葡萄、苜蓿、胡琴换了产地。商队是货的通道，也是观念的通道——经济和模因\textbf{咬合}的第一个大现场。
	\item \textbf{代价}：一条商路也是瘟疫的高速路。查士丁尼瘟疫沿贸易线传播——通道不筛善恶，只筛可达。
\end{itemize}

\textbf{用无你检验算一笔}：「拿掉丝绸之路，佛教会进中国吗？」——还能进（走海路、走使团），但概率显著降低。商路的贡献是正的，且不小；但不是「唯一原因」。这一笔「各占多少」，就是第 9 章的分账公式套在历史上的样子。

\begin{reader_anticipation}
	\item 价格有了「记分牌」，那谁在记分、能不能作弊？→ 记分的是千万次「贵了不买」的校准；作弊的路子第 14 章会看到——规则加稀缺，必有套利。
	\item 丝绸之路的「各占多少」真能算出来吗？→ 区间能算，精确值不能——历史案例的正确姿势是给区间、标把握，第 15 章失败清单第 5 条讲为什么。
\end{reader_anticipation}

\paragraph*{逻辑衔接}
	经济通道的账记完了：稀缺定规则，价格是记分牌，预期会放大，商路接文明。但「等价交换」这个词你信了很多年——下一章拆它：交换里有些东西\textbf{不许标价}，一标价，政治就上场了。

\begin{thinkbox}
\begin{enumerate}
	\item 用「摊主动机 $\rightarrow$ 全城菜价」的三层，拆一个你最近疑惑的价格：奶茶为什么涨价、演唱会门票为什么秒光、你行业的工资为什么涨/跌。约束变了哪一个？
	\item \textbf{反事实}：如果古代丝绸之路上「过路费」高到无利可图——贸易被压死——佛教还会在唐朝之前进中国吗？账要分两笔：经济这口锅占几成，模因那口锅占几成？
	\item \textbf{反事实}：如果猪周期里所有养猪户\textbf{同时}看得到全国存栏数据（信息完全透明），蛛网会消失吗？预期还在，振幅会变大还是变小？
\end{enumerate}
\end{thinkbox}

\end{document}
